\documentclass[11pt]{article}

\usepackage[utf8]{inputenc}
\usepackage[T1]{fontenc}
\usepackage{geometry}
\usepackage{amsmath}
\usepackage{amssymb}
\usepackage{booktabs}
\usepackage{hyperref}
\usepackage{xurl}
\usepackage{xcolor}
\usepackage{listings}
\usepackage{enumitem}
\usepackage{parskip}
\usepackage{graphicx}
\graphicspath{{../images/}}

\geometry{margin=1in}

\lstset{
  language=Java,
  basicstyle=\ttfamily\small,
  keywordstyle=\color{blue},
  commentstyle=\color{gray},
  stringstyle=\color{teal},
  showstringspaces=false,
  columns=fullflexible,
  frame=single,
  framerule=0.2pt,
  breaklines=true
}

\setlist[itemize]{topsep=0.3em,itemsep=0.2em}
\setlist[enumerate]{topsep=0.3em,itemsep=0.2em}

% Simple per-section source line.
\newcommand{\notesources}[1]{\par\noindent{\small\color{gray}\textit{Sources:} \sloppy #1\par}}

% Unit-wise source strings for this subject (used by slides/notes).
% Path is relative to the lecture `latex/` folder (the compilation CWD).
% OOPS with Java (CSEG1044) — unit-wise sources
%
% These are short, student-facing reference strings intended to be used with:
%   - \slidesources{...}  (in beamer slides)
%   - \notesources{...}   (in lecture notes)
%
% Style inspiration: other subjects in My-Subjects/ use semicolon-separated sources
% and include an "accessed" date for web documentation.

\newcommand{\OOPSUnitISources}{Schildt, \textit{Java: A Beginner's Guide} (10e), McGraw-Hill (Java fundamentals + OOP basics).; Oracle Java Tutorials: Getting Started + Language Basics + Classes and Objects (accessed \today).; Java SE API docs (e.g., \texttt{String}, \texttt{Scanner}) (accessed \today).}

\newcommand{\OOPSUnitIISources}{Schildt, \textit{Java: The Complete Reference} (12e), McGraw-Hill (classes, inheritance, packages, interfaces).; Bloch, \textit{Effective Java} (3e), Addison-Wesley, 2018 (design choices; interfaces vs abstract classes).; Oracle Java Tutorials: Classes and Objects + Interfaces + Packages (accessed \today).; Java SE API docs (e.g., \texttt{Comparable}, \texttt{Runnable}) (accessed \today).}

\newcommand{\OOPSUnitIIISources}{Schildt, \textit{Java: The Complete Reference} (12e) (nested/inner classes, exceptions, threads, I/O).; Oracle Java Tutorials: Nested Classes + Exceptions + Concurrency + Basic I/O (accessed \today).; Java SE API docs (e.g., \texttt{Thread}, \texttt{AutoCloseable}) (accessed \today).}

\newcommand{\OOPSUnitIVSources}{Schildt, \textit{Java: The Complete Reference} (12e) (generics, lambdas, Swing, JDBC).; Bloch, \textit{Effective Java} (3e) (generics + lambdas).; Oracle Java Tutorials: Generics + Lambda Expressions + Swing + JDBC Basics (accessed \today).; Java SE API docs (e.g., \texttt{java.util.function}, \texttt{javax.swing}, \texttt{java.sql}) (accessed \today).}

\newcommand{\OOPSUnitVSources}{Schildt, \textit{Java: The Complete Reference} (12e) (Collections Framework + wrapper classes).; Oracle Java docs: Collections Framework overview (accessed \today).; Java SE API docs (\texttt{java.util} collections; wrapper classes in \texttt{java.lang}) (accessed \today).}

\newcommand{\OOPSUnitVISources}{Git SCM: \textit{Pro Git} (2e) (version control workflow), accessed \today.; GitHub Docs (repositories, issues, pull requests), accessed \today.; Oracle Java Tutorials: Swing + JDBC (accessed \today).; JUnit 5 User Guide (accessed \today).; Apache Maven Project: Getting Started (accessed \today).; Gradle User Manual: Getting Started (accessed \today).; UML class diagrams: UML 2.x overview/spec (accessed \today).}

\newcommand{\OOPSMiscWhyInterfacesSources}{Schildt, \textit{Java: The Complete Reference} (12e) (interfaces).; Bloch, \textit{Effective Java} (3e) (prefer interfaces where appropriate).; Oracle Java Tutorials: Interfaces (accessed \today).; Java SE API docs (e.g., \texttt{Comparable}, \texttt{Runnable}) (accessed \today).}



\title{Object Oriented Programming (CSEG1044)\\Unit 02 -- Lecture 05 Notes\\Abstract Classes, `final`, and Interfaces (Design Choices)}
\author{Tofik Ali}
\date{\today}

\begin{document}
\maketitle
\tableofcontents

\section{Lecture Overview}
By now you have seen basic classes, inheritance, and an interface intro.
This lecture answers the practical design questions:
\begin{itemize}
  \item When should I use an \textbf{abstract class}?
  \item When should I use an \textbf{interface}?
  \item When should something be \textbf{final}?
\end{itemize}
\notesources{\OOPSUnitIISources}

\section{Syllabus Alignment (Unit 02)}
\textbf{Unit 02:} abstract classes, \texttt{final}, interfaces (built-in + user-defined)

\section{Learning Outcomes}
After this lecture, you should be able to:
\begin{itemize}
  \item Create an abstract class and implement abstract methods in subclasses.
  \item Explain the design intent of \texttt{final} variables/methods/classes.
  \item Write and use interfaces, including built-in \texttt{Comparable} and \texttt{Runnable}.
\end{itemize}

\section{Key Terms (Quick)}
\begin{itemize}
  \item \textbf{Abstract class}: class that cannot be instantiated; may include abstract methods.
  \item \textbf{Abstract method}: declared without body; must be implemented by concrete subclasses.
  \item \textbf{Interface}: contract specifying methods a class must provide.
  \item \textbf{Final}: ``cannot change'' (variable), ``cannot override'' (method), ``cannot extend'' (class).
  \item \textbf{Contract}: what is promised by a type (what callers can rely on).
\end{itemize}

\section{Detailed Notes}
\subsection{Abstract classes and methods}
Use an abstract class when:
\begin{itemize}
  \item you want to enforce a common API (method signatures),
  \item and you also want to share some implementation code across subclasses.
\end{itemize}

\textbf{Key properties:}
\begin{itemize}
  \item Cannot instantiate: \texttt{new Shape()} is illegal if \texttt{Shape} is abstract.
  \item Subclasses must implement all abstract methods (unless they are also abstract).
  \item An abstract class can contain concrete fields and methods.
\end{itemize}

\subsection{\texttt{final} keyword (design intent)}
\textbf{Final variable:} assigned once.
\begin{itemize}
  \item For primitives: value cannot change.
  \item For references: reference cannot change, but the object may still change unless the object is immutable.
\end{itemize}

\textbf{Final method:} cannot be overridden.
\begin{itemize}
  \item Used to lock down behavior and prevent subclasses from changing critical logic.
\end{itemize}

\textbf{Final class:} cannot be extended.
\begin{itemize}
  \item Example: \texttt{String} is final (one reason is security/immutability design).
\end{itemize}

\subsection{Interfaces: built-in + user-defined}
Use an interface when:
\begin{itemize}
  \item you want multiple unrelated classes to share a capability (e.g., \texttt{Comparable}),
  \item you want multiple implementations (e.g., different payment methods),
  \item you want flexible designs with low coupling.
\end{itemize}

\textbf{Built-in examples:}
\begin{itemize}
  \item \texttt{Comparable<T>}: defines natural ordering (\texttt{compareTo}). Used by sorting.
  \item \texttt{Runnable}: defines a task that can run in a thread (\texttt{run}).
\end{itemize}
\notesources{\OOPSUnitIISources}

\section{Worked Example: \texttt{Comparable} for sorting}
\begin{lstlisting}
import java.util.Arrays;

class Student implements Comparable<Student> {
    private final String rollNo;
    private final double cgpa;

    Student(String rollNo, double cgpa) {
        this.rollNo = rollNo;
        this.cgpa = cgpa;
    }

    @Override
    public int compareTo(Student other) {
        // descending order by cgpa
        return Double.compare(other.cgpa, this.cgpa);
    }

    @Override
    public String toString() {
        return rollNo + ":" + cgpa;
    }
}

public class Demo {
    public static void main(String[] args) {
        Student[] a = {
            new Student("21CS1001", 8.2),
            new Student("21CS1002", 9.1),
            new Student("21CS1003", 7.6)
        };
        Arrays.sort(a); // uses compareTo
        System.out.println(Arrays.toString(a));
    }
}
\end{lstlisting}

\section{In-Class Exercises (with brief solutions)}
\begin{enumerate}
  \item Can we create an object of an abstract class?\; \textbf{Answer:} No.
  \item What does \texttt{final class} mean?\; \textbf{Answer:} cannot be extended.
  \item Why do we use interfaces?\; \textbf{Answer:} contracts + low coupling + multiple implementations.
\end{enumerate}

\section{Self-Study Practice (no solutions)}
\begin{enumerate}
  \item Create an abstract class \texttt{Shape} with abstract method \texttt{area()} and implement \texttt{Circle} and \texttt{Rectangle}.
  \item Create an interface \texttt{Payable} and implement it for \texttt{Invoice} and \texttt{SalarySlip}.
  \item Implement \texttt{Runnable} in a class \texttt{HelloTask} and print a message in \texttt{run()} (we will start threads in Unit 03).
\end{enumerate}

\section{Checkpoint Questions (with sample answers)}
\textbf{Q1.} When do you prefer an interface over an abstract class?\par
\textbf{Sample answer:} Prefer an interface when you want a contract that many unrelated classes can implement and you want flexibility and low coupling. Use abstract class when you also want to share code and common state.

\vspace{0.6em}
\textbf{Q2.} What is the meaning of \texttt{final} in different contexts?\par
\textbf{Sample answer:} A final variable cannot be reassigned, a final method cannot be overridden, and a final class cannot be extended.

\end{document}
